\chapter{Tests and Doctests}
\label{ch:tests}

\standfirst{Two oracles, both runnable from the command line without a window,
and a ratchet that will not let a score fall or quietly rise.}

\section{SUnit}

The original of every xUnit, and it is here. Subclass \ct{TestCase}, write
methods whose names begin with \ct{test}, and assert:

\begin{st}
TestCase subclass: #AccountTest
    instanceVariableNames: 'account'
    classVariableNames: ''
    poolDictionaries: ''
    category: 'MyStuff-Tests'

AccountTest >> setUp
    account := Account new

AccountTest >> testStartsEmpty
    self assert: account balance equals: 0

AccountTest >> testDeposit
    account deposit: 100.
    self assert: account balance equals: 100

AccountTest >> testOverdraftIsRefused
    self should: [account withdraw: 1] raise: Error
\end{st}

\subsection{The assertions}

\begin{center}
\small
\begin{tabular}{@{}ll@{}}
\toprule
\ctp{assert:} \ctp{deny:} & a boolean, or a block answering one\\
\ctp{assert:description:} \ctp{deny:description:} & with a message\\
\ctp{assert:equals:} \ctp{deny:equals:} & reports both values on failure\\
\ctp{assert:identicalTo:} & \ctp{==}\\
\ctp{assert:closeTo:} \ctp{assert:closeTo:precision:} & floats\\
\ctp{assertEmpty:} \ctp{denyEmpty:} & \\
\ctp{assertCollection:equals:} & \\
\ctp{assertCollection:hasSameElements:} & order-independent\\
\ctp{should:raise:} \ctp{shouldnt:raise:} & exceptions\\
\bottomrule
\end{tabular}
\end{center}

Prefer \ct{assert:equals:} over \ct{assert:} with an \ct{=} inside it: when it
fails, the first tells you what the value actually was.

\subsection{setUp and tearDown}

\ct{setUp} runs before each test method and \ct{tearDown} after, so every test
starts from the same state. A fresh instance of the test class is made for
each test method, so instance variables cannot leak between tests.

\subsection{Running them}

From the command line, over a whole profile:

\begin{sh}
./st80 -bootstrap -profile profiles/mystuff.profile -tests
\end{sh}

\begin{plain}
12 run, 12 passed, 0 failed, 0 errors
\end{plain}

It exits non-zero if anything did not pass, so it drops straight into a
build script. In the image:

\begin{st}
AccountTest run
AccountTest suite run
(AccountTest selector: #testDeposit) run
\end{st}

\subsection{Failures and errors are different}

A \emph{failure} is an assertion that did not hold---the code did what it does
and you expected something else. An \emph{error} is an exception nobody
expected. \ct{TestResult} counts them separately, and the distinction is worth
keeping: a wall of errors usually means one thing is broken, and a wall of
failures usually means the code changed.

\begin{stnote}
\ct{TestFailure} \emph{is} an \ct{Error}, so the two handlers in
\ct{TestResult} are nested rather than side by side---the more specific one
has to be nearer up the sender chain, or every failure would be counted as an
error and the distinction would quietly stop working. That is a general lesson
about \ct{on:do:} and not just about SUnit.
\end{stnote}

\section{Doctests}

Pharo documents a method by putting examples in its comment, in a form meant
to be machine-readable:

\begin{st}
asBag
	"{1. 2} asBag = {2. 1} asBag >>> true"
	"#() asBag = Bag new >>> true"
	^Bag withAll: self
\end{st}

There are about 1,500 of them in the imported sources. They were written by
the people who wrote the methods, they cannot drift from the code they sit in,
and running them answers the question a port actually cares about---not "does
this source parse here" but "does it \emph{mean} here what it means there".

\begin{sh}
./st80 -bootstrap -profile profiles/st2026.profile -doctests lib/MyStuff/Account.class.st
\end{sh}

\begin{plain}
st80: 10 doctests in 3 methods of 1 files: 10 passed, 0 wrong,
      0 need something not here
\end{plain}

A doctest is checked by evaluating \ct{(expression) = (expected)}. Three
outcomes are told apart, because they are three different pieces of news:

\begin{center}
\small
\begin{tabular}{@{}lp{0.58\textwidth}@{}}
\toprule
passed & the expression answered the expected value\\
wrong & the method exists here and does something else. \textbf{This is a
  bug}\\
need something not here & it could not run: a class or selector this image has
  not got\\
\bottomrule
\end{tabular}
\end{center}

\begin{stnote}
The file being scanned need not be part of the image. \ct{-doctests} reads
examples out of any source file and runs them against whatever the profile
built---which makes it a general-purpose way to check a list of expressions.
This book's own examples are one such file; see below.
\end{stnote}

Doctests run under a much smaller bytecode budget than \ct{-eval}, on the
principle that an example written to show what a method does is not a
computation. One that does not answer quickly has hit something that does not
terminate.

\section{The ratchet}

\ct{make test} runs the unit suites and then every profile's own SUnit suites,
and holds the scores against \ct{tests/profiles.expected}:

\begin{plain}
# profile               run  passed
st2026                  12   12
pharo-announcements     43   43
pharo-time              633  633
pharo-weak              32   32
pharo-collections       469  469
\end{plain}

Both directions fail:

\begin{center}
\small
\begin{tabular}{@{}lp{0.60\textwidth}@{}}
\toprule
fewer passing & a regression --- the obvious one\\
fewer \emph{run} & tests that stopped being found. This looks like success and
  is worse\\
more of either & an unrecorded improvement, which nothing is protecting ---
  edit the file in the same commit that earns it\\
\bottomrule
\end{tabular}
\end{center}

The second row is the one that file exists for. A package whose tests stop
being found reports a perfect score, and the Chronology suites spent weeks
reporting nothing at all because four classes deadlocked the runner and the
number nobody printed was the number nobody read.

\section{Checking the book itself}

Everything in this tutorial that shows an answer is also a doctest, checked in
as \ct{tutorial/examples/TutorialExamples.class.st}:

\begin{sh}
cd tutorial && make verify
\end{sh}

If the book and the system ever disagree, that command is the referee. It is
the same mechanism the imported packages are held to, pointed at prose instead
of at source.
